\section{The settlement genus}

This section formalizes the site/settlement genus/species distinction: a bare
carrier/rule/topology (a site) paired with actual accreted content admitted by
its own law (a settlement). One abstract cell, instantiated independently at
material tiling, lenticular beadwork, causal boundary theory, and the
bitopological settling site's own \v{C}ech gluing --- none of the four built
with the others in mind.

\begin{definition}[Settlement]\label{def:settlement}\lean{Obsmetrologia.Settlement}\leanok
A settlement of a base is a pair of admitted content and a proof it is
admitted by the base's own law: \texttt{content : Content base} together
with \texttt{admits : Admits base content}.
\end{definition}

\begin{theorem}[Tiles are settlements of a finite local rule]\label{thm:settlement-of-tiles}\lean{Obsmetrologia.Settlement.ofTiles}\leanok\uses{def:settlement, thm:glass-substance-bundle}
A \texttt{Tiles} value is a settlement of its carrier/rule pair: the listed
pattern is the content, and the SFT \texttt{Path} admissibility is the
admission law.
\end{theorem}

\begin{theorem}[Beads are settlements of the lenticular hole-transition rule]\label{thm:settlement-of-beads-thread}\lean{Obsmetrologia.Settlement.ofBeadsThread}\leanok\uses{def:settlement, thm:glass-substance-bundle}
A \texttt{Beads} thread is a settlement of its lenticule carrier: the routed
thread is the content, and \texttt{ThreadLegal} --- every pass threading
through its lenticule via the hole-transition rule shared with the weave
crossing --- is the admission law.
\end{theorem}

\begin{theorem}[A Hole's effect relation is a settlement of its causal boundary substrate]\label{thm:settlement-of-hole}\lean{Obsmetrologia.Settlement.ofHole}\leanok\uses{def:settlement, thm:hole-shroud}
A \texttt{Hole}'s effect relation is a settlement of its causal substrate and
chosen readout: the effect relation is the content, and shroud-support
(\texttt{effect\_source\_shrouded}) is the admission law. The substrate and
readout alone already determine the shroud combinatorially, with no observed
activity presupposed --- the bare corner of this instance.
\end{theorem}

\begin{theorem}[A \v{C}ech datum's glued section is a settlement of the bitopological settling site]\label{thm:settlement-of-cech1}\lean{Obsmetrologia.Settlement.ofCech1}\leanok\uses{def:settlement, thm:glue, thm:cocycle, thm:h1-vanishes}
For an agreeing family (a \texttt{Cocycle}) over a \v{C}ech datum, the glued
global section is a settlement of the datum: the section is the content, and
being a \texttt{Coboundary} of the family is the admission law --- reusing
\texttt{glue}/\texttt{glue\_coboundary} rather than re-deriving gluing.
\end{theorem}

\begin{theorem}[The recurring settlement cell, met four times]\label{thm:settlement-recurring-cell}\lean{Obsmetrologia.Settlement.tiles_beads_hole_cech1_are_settlements}\leanok\uses{thm:settlement-of-tiles, thm:settlement-of-beads-thread, thm:settlement-of-hole, thm:settlement-of-cech1}
\texttt{Tiles}, a \texttt{Beads} thread, a \texttt{Hole}'s effect relation,
and a \v{C}ech datum's glued section are each a settlement of their own base
--- the same two-field shape, independently instantiated in four domains
built without each other in mind.
\end{theorem}
